% !TEX program = xelatex
\documentclass[a4paper,11pt,openany,fontset=none]{ctexrep}

\newcommand{\reporttitle}{稳健多资产配置策略}
\newcommand{\reportsubtitle}{20万元本金的条件性风险预算与复利框架}
\newcommand{\reportkicker}{A股多资产战略配置研究}
\newcommand{\reportscope}{流动性 | 利率债 | 红利权益ETF | 非杠杆黄金ETF | 研究用途}
\newcommand{\reportdate}{2026年8月12日}
\newcommand{\reportdatacutoff}{规则和产品资料截至2026年8月11日；模型标的与实施顺序于2026年8月12日明确；当前范围不使用个股价格或历史股息率快照}
\newcommand{\reporttype}{完整策略配置报告}
\newcommand{\reportauthor}{AStock Research}
\newcommand{\reporthouseview}{固定20万元本金的条件性模型目标为：CNY70,000的511010、CNY30,000的510880、通过穿透核验后才可启用的CNY30,000的512890，以及CNY20,000的518880；另保留CNY40,000受保障流动性和CNY10,000交易现金。当前ETF资格包未运行，直接股票和REITs初始均为0\%。}
\newcommand{\reportquality}{当前策略主张使用已本地归档的规则、交易所和产品原始资料；不使用个股价格或历史股息率快照。当前价格、IOPV、折溢价、价差、流动性、账户资格和合并持仓均须在实施日重新核验。}
\newcommand{\reportdisclaimer}{本报告基于公开资料整理，仅作研究与教育参考，不构成证券买卖建议，不保证收益或本金安全。}

% Reusable institutional-style preamble for XeLaTeX AStock research reports.
% Usage: define \reporttitle and \reportdate, then % Reusable institutional-style preamble for XeLaTeX AStock research reports.
% Usage: define \reporttitle and \reportdate, then % Reusable institutional-style preamble for XeLaTeX AStock research reports.
% Usage: define \reporttitle and \reportdate, then \input{.agents/templates/preamble.tex}.
% Design benchmark: JPM first-page dashboard, BlackRock house-view framing,
% Vanguard scenario discipline, and AB/JPM chartbook exhibit structure.

\providecommand{\reporttitle}{AStock Research Report}
\providecommand{\reportsubtitle}{House View, Evidence, Valuation and Risk}
\providecommand{\reportkicker}{INSTITUTIONAL EQUITY RESEARCH}
\providecommand{\reportscope}{CHINA A-SHARES}
\providecommand{\reportdate}{\today}
\providecommand{\reportdatacutoff}{Data cutoff not specified}
\providecommand{\reporttype}{Research note}
\providecommand{\reportauthor}{AStock Research Agent}
\providecommand{\reporthouseview}{State the house view in one decision-useful paragraph.}
\providecommand{\reportquality}{Source quality not specified.}
\providecommand{\reportdisclaimer}{本报告基于公开资料整理，不构成任何证券买卖建议。}

% ============ 页面布局 ============
\usepackage[top=1.8cm, bottom=1.9cm, left=1.8cm, right=1.8cm]{geometry}
\setlength{\parskip}{0.3em}
\setlength{\parindent}{2em}
\renewcommand{\baselinestretch}{1.12}
\raggedbottom

% ============ Fonts ============
% Prefer the installed Noto CJK family when available (for Linux builds).
% Otherwise retain the original MacTeX typography. This avoids platform
% probing via shell escape and keeps the source portable across both engines.
\IfFontExistsTF{Noto Serif CJK SC}{
  \setCJKmainfont{Noto Serif CJK SC}
  \setCJKsansfont{Noto Sans CJK SC}
  \setCJKmonofont{Noto Sans CJK SC}
  \setmainfont{Liberation Serif}
  \newcommand{\kaishu}{\CJKfontspec{Noto Serif CJK SC}}
  \newcommand{\heiti}{\CJKfontspec{Noto Sans CJK SC}}
  \newcommand{\songti}{\CJKfontspec{Noto Serif CJK SC}}
}{
  \setCJKmainfont[BoldFont={Heiti SC}]{STSong}
  \setCJKsansfont{Heiti SC}
  \setCJKmonofont{Heiti SC}
  \setmainfont{Times New Roman}
  \newcommand{\kaishu}{\CJKfontspec{STKaiti}}
  \newcommand{\heiti}{\CJKfontspec{Heiti SC}}
  \newcommand{\songti}{\CJKfontspec{STSong}}
}

% ============ 颜色（IB Navy/Grey palette） ============
\usepackage{xcolor}
\definecolor{navy}{HTML}{003366}
\definecolor{deepnavy}{HTML}{0B1F33}
\definecolor{steelgrey}{HTML}{4A5568}
\definecolor{mutedgrey}{HTML}{718096}
\definecolor{lightgrey}{HTML}{F7FAFC}
\definecolor{panelgrey}{HTML}{EDF2F7}
\definecolor{riskred}{HTML}{C53030}
\definecolor{riskamber}{HTML}{D69E2E}
\definecolor{riskgreen}{HTML}{38A169}
\definecolor{riskblue}{HTML}{2B6CB0}
\definecolor{accentblue}{HTML}{2B6CB0}
\definecolor{nvgreen}{HTML}{76B900}
\definecolor{rulegrey}{HTML}{CBD5E0}

% ============ 表格 ============
\usepackage{booktabs}
\usepackage{longtable}
\usepackage{tabularx}
\usepackage{multirow}
\usepackage{array}
\usepackage{makecell}
\usepackage{colortbl}
\renewcommand{\arraystretch}{1.12}
\newcolumntype{Y}{>{\centering\arraybackslash}X}
\newcolumntype{L}[1]{>{\raggedright\arraybackslash}p{#1}}
\newcolumntype{C}[1]{>{\centering\arraybackslash}p{#1}}
\newcolumntype{R}[1]{>{\raggedleft\arraybackslash}p{#1}}

% ============ 图形 ============
\usepackage{graphicx}
\usepackage{tikz}
\usetikzlibrary{shapes,arrows.meta,positioning,calc,backgrounds,fit,decorations.pathreplacing}
\usepackage{pgfplots}
\usepgfplotslibrary{polar}
\pgfplotsset{compat=1.18}
\usepackage[edges]{forest}

% ============ 页眉页脚（报告标题和日期需在main.tex中覆盖） ============
\usepackage{fancyhdr}
\pagestyle{fancy}
\fancyhf{}
\fancyhead[L]{\scriptsize\color{steelgrey} \reporttitle}
\fancyhead[R]{\scriptsize\color{steelgrey} \reportdate}
\fancyfoot[C]{\small\color{steelgrey} \thepage}
\fancyfoot[R]{\tiny\color{steelgrey} \reportdisclaimer}
\renewcommand{\headrulewidth}{0.4pt}
\renewcommand{\footrulewidth}{0.2pt}

% ============ 章节格式 ============
\usepackage{titlesec}
\titleformat{\chapter}[block]
  {\normalfont\LARGE\bfseries\color{navy}}
  {\thechapter}{0.9em}{}
\titlespacing*{\chapter}{0pt}{0.4em}{1.0em}
\titleformat{\section}
  {\normalfont\large\bfseries\color{navy}}
  {\thesection}{1em}{}
\titleformat{\subsection}
  {\normalfont\normalsize\bfseries\color{steelgrey}}
  {\thesubsection}{1em}{}

% ============ 超链接 ============
\usepackage{hyperref}
\hypersetup{
  colorlinks=true,
  linkcolor=navy,
  urlcolor=navy,
  citecolor=navy,
}

% ============ 其他工具 ============
\usepackage{enumitem}
\setlist{nosep, leftmargin=2em}
\IfFileExists{fontawesome5.sty}{\usepackage{fontawesome5}}{}
% ============ Class B 全局兜底（xeCJK + 缺字 fallback 阻断）：
% 当 fontawesome5 因 Times New Roman 缺字 / \FA 命名空间失败时，
% 禁止输出 aExclamationTriangle / aCloud 裸文本，一律降级为 amssymb 原生符号。
% 注：amssymb 在 L112 已加载，此处不重复。所有 \def 在本段落。
\def\faExclamationTriangle{\ensuremath{\textcolor{riskred}{\blacktriangle}}\kern-0.25em}
\def\faCloud{\ensuremath{\clubsuit}\kern-0.1em}
\def\faYenSign{\ensuremath{\yen}}
\def\faIndustry{\ensuremath{\S}}
\def\faCogs{\ensuremath{\ast}}
\def\faFlagCheckered{\ensuremath{\blacktriangleright}}
\def\faMemory{\textbf{[RAM]}}
\def\faServer{\textbf{[DC]}}
\def\faGlobe{\ensuremath{\diamond}}
\def\faStar{\ensuremath{\bigstar}}
\def\faMicrochip{\textbf{[$\mu$P]}}
\def\faLayerGroup{\textbf{[LAY]}}
\def\faCubes{\ensuremath{\boxtimes}}
\usepackage{amssymb}
\usepackage{pifont}
\usepackage{tcolorbox}
\tcbuselibrary{skins,breakable}
\usepackage{caption}
\captionsetup{font=small, skip=4pt}
\usepackage{float}
\usepackage{needspace}
\setlength{\textfloatsep}{8pt plus 2pt minus 2pt}
\setlength{\floatsep}{6pt plus 2pt minus 2pt}
\setlength{\intextsep}{6pt plus 2pt minus 2pt}
\setlength{\abovecaptionskip}{4pt}
\setlength{\belowcaptionskip}{2pt}

% ============ 自定义命令 ============
\newcommand{\rmark}{\textcolor{riskred}{\ding{108}}}
\newcommand{\amark}{\textcolor{riskamber}{\ding{108}}}
\newcommand{\gmark}{\textcolor{riskgreen}{\ding{108}}}
\newcommand{\starfive}{$\bigstar\bigstar\bigstar\bigstar\bigstar$}
\newcommand{\starfour}{$\bigstar\bigstar\bigstar\bigstar$}
\newcommand{\starthree}{$\bigstar\bigstar\bigstar$}
% 资料来源注：使用楷体（STKaiti）而非斜体，避免中文无斜体字形导致的字体回退警告
% \mdseries 避免继承周围 \bfseries 导致 STKaiti 粗体缺失警告
% \sloppy + \hbadness 允许长文件名/URL 在必要时断行
% 局部重定义 \textbf → \heiti：STKaiti 无粗体，用黑体替代强调
\newcommand{\sourcenote}[1]{\par{\scriptsize\mdseries\color{mutedgrey}\kaishu\sloppy\hbadness=9999\def\textbf##1{{\heiti ##1}} 资料来源：#1}\vspace{0.2em}}
\newcommand{\qualitynote}[1]{\par{\scriptsize\mdseries\color{mutedgrey}\sloppy\def\textbf##1{{\heiti ##1}} 数据质量：#1}\vspace{0.2em}}
\newcommand{\evidencetag}[1]{%
  \tikz[baseline=(x.base)]\node[
    xshift=0pt,
    inner xsep=3.5pt,
    inner ysep=1.2pt,
    rounded corners=1.5pt,
    fill=panelgrey,
    draw=rulegrey,
    text=steelgrey,
    font=\scriptsize
  ] (x) {#1};%
}
\newcommand{\stance}[2]{\textbf{\textcolor{#1}{#2}}}

% 信息框
\newtcolorbox{keyinsight}[1][]{
  colback=lightgrey, colframe=navy,
  fonttitle=\bfseries, title={#1}, breakable,
  arc=2pt, boxrule=0.5pt, left=8pt, right=8pt, top=6pt, bottom=6pt
}
\newtcolorbox{riskbox}[1][]{
  colback=red!5, colframe=riskred,
  fonttitle=\bfseries, title={#1}, breakable,
  arc=2pt, boxrule=0.5pt, left=8pt, right=8pt, top=6pt, bottom=6pt
}
\newtcolorbox{dashboardbox}[1][]{
  enhanced, breakable,
  colback=white, colframe=navy,
  boxrule=0.7pt, arc=2pt,
  fonttitle=\bfseries\color{white},
  coltitle=white, colbacktitle=navy,
  title={#1},
  left=9pt, right=9pt, top=7pt, bottom=7pt
}
\newtcolorbox{houseviewbox}[1][]{
  enhanced, breakable,
  colback=lightgrey, colframe=deepnavy,
  borderline west={3pt}{0pt}{accentblue},
  boxrule=0.4pt, arc=2pt,
  fonttitle=\bfseries, title={#1},
  left=9pt, right=9pt, top=7pt, bottom=7pt
}
\newtcolorbox{sourcequalitybox}[1][]{
  enhanced, breakable,
  colback=yellow!6, colframe=riskamber,
  boxrule=0.45pt, arc=2pt,
  fonttitle=\bfseries, title={#1},
  left=9pt, right=9pt, top=7pt, bottom=7pt
}
% exhibitbox：表/图外框。left=7pt+right=7pt=14pt 内边距，before upper 内的 \linewidth 已扣除该内边距。
% 约定：\exhibitboxwidth = 真实内边距宽（tcolorbox 已扣好 padding 的 \linewidth）。
% 由于 tabularx/longtable 宽度参数的语义是「外部总宽（含两侧 \tabcolsep）」，内部列实际可用 = W - 2\tabcolsep。
% 所以 tabularx 直接传 \exhibitboxwidth（总宽 == 内边距宽），而其他用途（resizebox、tikz、p{...}）用 \exhibitwidth = \exhibitboxwidth - 2\tabcolsep。
\newlength{\exhibitboxwidth}
\newlength{\exhibitwidth}
\newtcolorbox{exhibitbox}[1][]{
  enhanced, breakable,
  colback=white, colframe=rulegrey,
  borderline north={1.2pt}{0pt}{navy},
  boxrule=0.35pt, arc=1pt,
  fonttitle=\bfseries\color{navy}, title={#1},
  left=7pt, right=7pt, top=6pt, bottom=6pt,
  before upper={\global\exhibitboxwidth=\linewidth\relax
    \global\exhibitwidth=\dimexpr\linewidth-2\tabcolsep\relax},
}
\newtcolorbox{disclosurebox}[1][]{
  enhanced, breakable,
  colback=lightgrey, colframe=steelgrey,
  boxrule=0.35pt, arc=1pt,
  fonttitle=\bfseries, title={#1},
  left=8pt, right=8pt, top=6pt, bottom=6pt
}

% ============ First-page dashboard cover ============
\newcommand{\astockcover}{%
\begin{titlepage}
\thispagestyle{empty}
\begin{tikzpicture}[remember picture, overlay]
  \fill[deepnavy] (current page.north west) rectangle ([yshift=-2.55cm]current page.north east);
  \fill[accentblue] ([yshift=-2.55cm]current page.north west) rectangle ([yshift=-2.66cm]current page.north east);
  \fill[deepnavy] (current page.south west) rectangle ([yshift=1.2cm]current page.south east);
\end{tikzpicture}
\vspace*{0.62cm}
{\color{white}\sffamily\bfseries\small \reportkicker}\\[0.18cm]
{\color{steelgrey}\sffamily\scriptsize \reportscope \quad|\quad \reporttype \quad|\quad \reportdate}\\[1.0cm]
{\fontsize{25}{30}\selectfont\bfseries\color{navy}\reporttitle}\\[0.25cm]
{\fontsize{13}{17}\selectfont\color{steelgrey}\reportsubtitle}\\[0.75cm]
\begin{dashboardbox}[决策摘要]
\begin{tabularx}{\textwidth}{L{3.0cm}X L{3.0cm}X}
\textbf{报告类型} & \reporttype & \textbf{数据截至} & \reportdatacutoff \\
\textbf{覆盖范围} & \reportscope & \textbf{作者} & \reportauthor \\
\end{tabularx}
\vspace{0.45em}
\begin{houseviewbox}[本机构观点]
\reporthouseview
\end{houseviewbox}
\begin{sourcequalitybox}[证据质量]
\reportquality
\end{sourcequalitybox}
\end{dashboardbox}
\vfill
{\scriptsize\color{steelgrey}\reportdisclaimer}
\end{titlepage}
}
.
% Design benchmark: JPM first-page dashboard, BlackRock house-view framing,
% Vanguard scenario discipline, and AB/JPM chartbook exhibit structure.

\providecommand{\reporttitle}{AStock Research Report}
\providecommand{\reportsubtitle}{House View, Evidence, Valuation and Risk}
\providecommand{\reportkicker}{INSTITUTIONAL EQUITY RESEARCH}
\providecommand{\reportscope}{CHINA A-SHARES}
\providecommand{\reportdate}{\today}
\providecommand{\reportdatacutoff}{Data cutoff not specified}
\providecommand{\reporttype}{Research note}
\providecommand{\reportauthor}{AStock Research Agent}
\providecommand{\reporthouseview}{State the house view in one decision-useful paragraph.}
\providecommand{\reportquality}{Source quality not specified.}
\providecommand{\reportdisclaimer}{本报告基于公开资料整理，不构成任何证券买卖建议。}

% ============ 页面布局 ============
\usepackage[top=1.8cm, bottom=1.9cm, left=1.8cm, right=1.8cm]{geometry}
\setlength{\parskip}{0.3em}
\setlength{\parindent}{2em}
\renewcommand{\baselinestretch}{1.12}
\raggedbottom

% ============ Fonts ============
% Prefer the installed Noto CJK family when available (for Linux builds).
% Otherwise retain the original MacTeX typography. This avoids platform
% probing via shell escape and keeps the source portable across both engines.
\IfFontExistsTF{Noto Serif CJK SC}{
  \setCJKmainfont{Noto Serif CJK SC}
  \setCJKsansfont{Noto Sans CJK SC}
  \setCJKmonofont{Noto Sans CJK SC}
  \setmainfont{Liberation Serif}
  \newcommand{\kaishu}{\CJKfontspec{Noto Serif CJK SC}}
  \newcommand{\heiti}{\CJKfontspec{Noto Sans CJK SC}}
  \newcommand{\songti}{\CJKfontspec{Noto Serif CJK SC}}
}{
  \setCJKmainfont[BoldFont={Heiti SC}]{STSong}
  \setCJKsansfont{Heiti SC}
  \setCJKmonofont{Heiti SC}
  \setmainfont{Times New Roman}
  \newcommand{\kaishu}{\CJKfontspec{STKaiti}}
  \newcommand{\heiti}{\CJKfontspec{Heiti SC}}
  \newcommand{\songti}{\CJKfontspec{STSong}}
}

% ============ 颜色（IB Navy/Grey palette） ============
\usepackage{xcolor}
\definecolor{navy}{HTML}{003366}
\definecolor{deepnavy}{HTML}{0B1F33}
\definecolor{steelgrey}{HTML}{4A5568}
\definecolor{mutedgrey}{HTML}{718096}
\definecolor{lightgrey}{HTML}{F7FAFC}
\definecolor{panelgrey}{HTML}{EDF2F7}
\definecolor{riskred}{HTML}{C53030}
\definecolor{riskamber}{HTML}{D69E2E}
\definecolor{riskgreen}{HTML}{38A169}
\definecolor{riskblue}{HTML}{2B6CB0}
\definecolor{accentblue}{HTML}{2B6CB0}
\definecolor{nvgreen}{HTML}{76B900}
\definecolor{rulegrey}{HTML}{CBD5E0}

% ============ 表格 ============
\usepackage{booktabs}
\usepackage{longtable}
\usepackage{tabularx}
\usepackage{multirow}
\usepackage{array}
\usepackage{makecell}
\usepackage{colortbl}
\renewcommand{\arraystretch}{1.12}
\newcolumntype{Y}{>{\centering\arraybackslash}X}
\newcolumntype{L}[1]{>{\raggedright\arraybackslash}p{#1}}
\newcolumntype{C}[1]{>{\centering\arraybackslash}p{#1}}
\newcolumntype{R}[1]{>{\raggedleft\arraybackslash}p{#1}}

% ============ 图形 ============
\usepackage{graphicx}
\usepackage{tikz}
\usetikzlibrary{shapes,arrows.meta,positioning,calc,backgrounds,fit,decorations.pathreplacing}
\usepackage{pgfplots}
\usepgfplotslibrary{polar}
\pgfplotsset{compat=1.18}
\usepackage[edges]{forest}

% ============ 页眉页脚（报告标题和日期需在main.tex中覆盖） ============
\usepackage{fancyhdr}
\pagestyle{fancy}
\fancyhf{}
\fancyhead[L]{\scriptsize\color{steelgrey} \reporttitle}
\fancyhead[R]{\scriptsize\color{steelgrey} \reportdate}
\fancyfoot[C]{\small\color{steelgrey} \thepage}
\fancyfoot[R]{\tiny\color{steelgrey} \reportdisclaimer}
\renewcommand{\headrulewidth}{0.4pt}
\renewcommand{\footrulewidth}{0.2pt}

% ============ 章节格式 ============
\usepackage{titlesec}
\titleformat{\chapter}[block]
  {\normalfont\LARGE\bfseries\color{navy}}
  {\thechapter}{0.9em}{}
\titlespacing*{\chapter}{0pt}{0.4em}{1.0em}
\titleformat{\section}
  {\normalfont\large\bfseries\color{navy}}
  {\thesection}{1em}{}
\titleformat{\subsection}
  {\normalfont\normalsize\bfseries\color{steelgrey}}
  {\thesubsection}{1em}{}

% ============ 超链接 ============
\usepackage{hyperref}
\hypersetup{
  colorlinks=true,
  linkcolor=navy,
  urlcolor=navy,
  citecolor=navy,
}

% ============ 其他工具 ============
\usepackage{enumitem}
\setlist{nosep, leftmargin=2em}
\IfFileExists{fontawesome5.sty}{\usepackage{fontawesome5}}{}
% ============ Class B 全局兜底（xeCJK + 缺字 fallback 阻断）：
% 当 fontawesome5 因 Times New Roman 缺字 / \FA 命名空间失败时，
% 禁止输出 aExclamationTriangle / aCloud 裸文本，一律降级为 amssymb 原生符号。
% 注：amssymb 在 L112 已加载，此处不重复。所有 \def 在本段落。
\def\faExclamationTriangle{\ensuremath{\textcolor{riskred}{\blacktriangle}}\kern-0.25em}
\def\faCloud{\ensuremath{\clubsuit}\kern-0.1em}
\def\faYenSign{\ensuremath{\yen}}
\def\faIndustry{\ensuremath{\S}}
\def\faCogs{\ensuremath{\ast}}
\def\faFlagCheckered{\ensuremath{\blacktriangleright}}
\def\faMemory{\textbf{[RAM]}}
\def\faServer{\textbf{[DC]}}
\def\faGlobe{\ensuremath{\diamond}}
\def\faStar{\ensuremath{\bigstar}}
\def\faMicrochip{\textbf{[$\mu$P]}}
\def\faLayerGroup{\textbf{[LAY]}}
\def\faCubes{\ensuremath{\boxtimes}}
\usepackage{amssymb}
\usepackage{pifont}
\usepackage{tcolorbox}
\tcbuselibrary{skins,breakable}
\usepackage{caption}
\captionsetup{font=small, skip=4pt}
\usepackage{float}
\usepackage{needspace}
\setlength{\textfloatsep}{8pt plus 2pt minus 2pt}
\setlength{\floatsep}{6pt plus 2pt minus 2pt}
\setlength{\intextsep}{6pt plus 2pt minus 2pt}
\setlength{\abovecaptionskip}{4pt}
\setlength{\belowcaptionskip}{2pt}

% ============ 自定义命令 ============
\newcommand{\rmark}{\textcolor{riskred}{\ding{108}}}
\newcommand{\amark}{\textcolor{riskamber}{\ding{108}}}
\newcommand{\gmark}{\textcolor{riskgreen}{\ding{108}}}
\newcommand{\starfive}{$\bigstar\bigstar\bigstar\bigstar\bigstar$}
\newcommand{\starfour}{$\bigstar\bigstar\bigstar\bigstar$}
\newcommand{\starthree}{$\bigstar\bigstar\bigstar$}
% 资料来源注：使用楷体（STKaiti）而非斜体，避免中文无斜体字形导致的字体回退警告
% \mdseries 避免继承周围 \bfseries 导致 STKaiti 粗体缺失警告
% \sloppy + \hbadness 允许长文件名/URL 在必要时断行
% 局部重定义 \textbf → \heiti：STKaiti 无粗体，用黑体替代强调
\newcommand{\sourcenote}[1]{\par{\scriptsize\mdseries\color{mutedgrey}\kaishu\sloppy\hbadness=9999\def\textbf##1{{\heiti ##1}} 资料来源：#1}\vspace{0.2em}}
\newcommand{\qualitynote}[1]{\par{\scriptsize\mdseries\color{mutedgrey}\sloppy\def\textbf##1{{\heiti ##1}} 数据质量：#1}\vspace{0.2em}}
\newcommand{\evidencetag}[1]{%
  \tikz[baseline=(x.base)]\node[
    xshift=0pt,
    inner xsep=3.5pt,
    inner ysep=1.2pt,
    rounded corners=1.5pt,
    fill=panelgrey,
    draw=rulegrey,
    text=steelgrey,
    font=\scriptsize
  ] (x) {#1};%
}
\newcommand{\stance}[2]{\textbf{\textcolor{#1}{#2}}}

% 信息框
\newtcolorbox{keyinsight}[1][]{
  colback=lightgrey, colframe=navy,
  fonttitle=\bfseries, title={#1}, breakable,
  arc=2pt, boxrule=0.5pt, left=8pt, right=8pt, top=6pt, bottom=6pt
}
\newtcolorbox{riskbox}[1][]{
  colback=red!5, colframe=riskred,
  fonttitle=\bfseries, title={#1}, breakable,
  arc=2pt, boxrule=0.5pt, left=8pt, right=8pt, top=6pt, bottom=6pt
}
\newtcolorbox{dashboardbox}[1][]{
  enhanced, breakable,
  colback=white, colframe=navy,
  boxrule=0.7pt, arc=2pt,
  fonttitle=\bfseries\color{white},
  coltitle=white, colbacktitle=navy,
  title={#1},
  left=9pt, right=9pt, top=7pt, bottom=7pt
}
\newtcolorbox{houseviewbox}[1][]{
  enhanced, breakable,
  colback=lightgrey, colframe=deepnavy,
  borderline west={3pt}{0pt}{accentblue},
  boxrule=0.4pt, arc=2pt,
  fonttitle=\bfseries, title={#1},
  left=9pt, right=9pt, top=7pt, bottom=7pt
}
\newtcolorbox{sourcequalitybox}[1][]{
  enhanced, breakable,
  colback=yellow!6, colframe=riskamber,
  boxrule=0.45pt, arc=2pt,
  fonttitle=\bfseries, title={#1},
  left=9pt, right=9pt, top=7pt, bottom=7pt
}
% exhibitbox：表/图外框。left=7pt+right=7pt=14pt 内边距，before upper 内的 \linewidth 已扣除该内边距。
% 约定：\exhibitboxwidth = 真实内边距宽（tcolorbox 已扣好 padding 的 \linewidth）。
% 由于 tabularx/longtable 宽度参数的语义是「外部总宽（含两侧 \tabcolsep）」，内部列实际可用 = W - 2\tabcolsep。
% 所以 tabularx 直接传 \exhibitboxwidth（总宽 == 内边距宽），而其他用途（resizebox、tikz、p{...}）用 \exhibitwidth = \exhibitboxwidth - 2\tabcolsep。
\newlength{\exhibitboxwidth}
\newlength{\exhibitwidth}
\newtcolorbox{exhibitbox}[1][]{
  enhanced, breakable,
  colback=white, colframe=rulegrey,
  borderline north={1.2pt}{0pt}{navy},
  boxrule=0.35pt, arc=1pt,
  fonttitle=\bfseries\color{navy}, title={#1},
  left=7pt, right=7pt, top=6pt, bottom=6pt,
  before upper={\global\exhibitboxwidth=\linewidth\relax
    \global\exhibitwidth=\dimexpr\linewidth-2\tabcolsep\relax},
}
\newtcolorbox{disclosurebox}[1][]{
  enhanced, breakable,
  colback=lightgrey, colframe=steelgrey,
  boxrule=0.35pt, arc=1pt,
  fonttitle=\bfseries, title={#1},
  left=8pt, right=8pt, top=6pt, bottom=6pt
}

% ============ First-page dashboard cover ============
\newcommand{\astockcover}{%
\begin{titlepage}
\thispagestyle{empty}
\begin{tikzpicture}[remember picture, overlay]
  \fill[deepnavy] (current page.north west) rectangle ([yshift=-2.55cm]current page.north east);
  \fill[accentblue] ([yshift=-2.55cm]current page.north west) rectangle ([yshift=-2.66cm]current page.north east);
  \fill[deepnavy] (current page.south west) rectangle ([yshift=1.2cm]current page.south east);
\end{tikzpicture}
\vspace*{0.62cm}
{\color{white}\sffamily\bfseries\small \reportkicker}\\[0.18cm]
{\color{steelgrey}\sffamily\scriptsize \reportscope \quad|\quad \reporttype \quad|\quad \reportdate}\\[1.0cm]
{\fontsize{25}{30}\selectfont\bfseries\color{navy}\reporttitle}\\[0.25cm]
{\fontsize{13}{17}\selectfont\color{steelgrey}\reportsubtitle}\\[0.75cm]
\begin{dashboardbox}[决策摘要]
\begin{tabularx}{\textwidth}{L{3.0cm}X L{3.0cm}X}
\textbf{报告类型} & \reporttype & \textbf{数据截至} & \reportdatacutoff \\
\textbf{覆盖范围} & \reportscope & \textbf{作者} & \reportauthor \\
\end{tabularx}
\vspace{0.45em}
\begin{houseviewbox}[本机构观点]
\reporthouseview
\end{houseviewbox}
\begin{sourcequalitybox}[证据质量]
\reportquality
\end{sourcequalitybox}
\end{dashboardbox}
\vfill
{\scriptsize\color{steelgrey}\reportdisclaimer}
\end{titlepage}
}
.
% Design benchmark: JPM first-page dashboard, BlackRock house-view framing,
% Vanguard scenario discipline, and AB/JPM chartbook exhibit structure.

\providecommand{\reporttitle}{AStock Research Report}
\providecommand{\reportsubtitle}{House View, Evidence, Valuation and Risk}
\providecommand{\reportkicker}{INSTITUTIONAL EQUITY RESEARCH}
\providecommand{\reportscope}{CHINA A-SHARES}
\providecommand{\reportdate}{\today}
\providecommand{\reportdatacutoff}{Data cutoff not specified}
\providecommand{\reporttype}{Research note}
\providecommand{\reportauthor}{AStock Research Agent}
\providecommand{\reporthouseview}{State the house view in one decision-useful paragraph.}
\providecommand{\reportquality}{Source quality not specified.}
\providecommand{\reportdisclaimer}{本报告基于公开资料整理，不构成任何证券买卖建议。}

% ============ 页面布局 ============
\usepackage[top=1.8cm, bottom=1.9cm, left=1.8cm, right=1.8cm]{geometry}
\setlength{\parskip}{0.3em}
\setlength{\parindent}{2em}
\renewcommand{\baselinestretch}{1.12}
\raggedbottom

% ============ Fonts ============
% Prefer the installed Noto CJK family when available (for Linux builds).
% Otherwise retain the original MacTeX typography. This avoids platform
% probing via shell escape and keeps the source portable across both engines.
\IfFontExistsTF{Noto Serif CJK SC}{
  \setCJKmainfont{Noto Serif CJK SC}
  \setCJKsansfont{Noto Sans CJK SC}
  \setCJKmonofont{Noto Sans CJK SC}
  \setmainfont{Liberation Serif}
  \newcommand{\kaishu}{\CJKfontspec{Noto Serif CJK SC}}
  \newcommand{\heiti}{\CJKfontspec{Noto Sans CJK SC}}
  \newcommand{\songti}{\CJKfontspec{Noto Serif CJK SC}}
}{
  \setCJKmainfont[BoldFont={Heiti SC}]{STSong}
  \setCJKsansfont{Heiti SC}
  \setCJKmonofont{Heiti SC}
  \setmainfont{Times New Roman}
  \newcommand{\kaishu}{\CJKfontspec{STKaiti}}
  \newcommand{\heiti}{\CJKfontspec{Heiti SC}}
  \newcommand{\songti}{\CJKfontspec{STSong}}
}

% ============ 颜色（IB Navy/Grey palette） ============
\usepackage{xcolor}
\definecolor{navy}{HTML}{003366}
\definecolor{deepnavy}{HTML}{0B1F33}
\definecolor{steelgrey}{HTML}{4A5568}
\definecolor{mutedgrey}{HTML}{718096}
\definecolor{lightgrey}{HTML}{F7FAFC}
\definecolor{panelgrey}{HTML}{EDF2F7}
\definecolor{riskred}{HTML}{C53030}
\definecolor{riskamber}{HTML}{D69E2E}
\definecolor{riskgreen}{HTML}{38A169}
\definecolor{riskblue}{HTML}{2B6CB0}
\definecolor{accentblue}{HTML}{2B6CB0}
\definecolor{nvgreen}{HTML}{76B900}
\definecolor{rulegrey}{HTML}{CBD5E0}

% ============ 表格 ============
\usepackage{booktabs}
\usepackage{longtable}
\usepackage{tabularx}
\usepackage{multirow}
\usepackage{array}
\usepackage{makecell}
\usepackage{colortbl}
\renewcommand{\arraystretch}{1.12}
\newcolumntype{Y}{>{\centering\arraybackslash}X}
\newcolumntype{L}[1]{>{\raggedright\arraybackslash}p{#1}}
\newcolumntype{C}[1]{>{\centering\arraybackslash}p{#1}}
\newcolumntype{R}[1]{>{\raggedleft\arraybackslash}p{#1}}

% ============ 图形 ============
\usepackage{graphicx}
\usepackage{tikz}
\usetikzlibrary{shapes,arrows.meta,positioning,calc,backgrounds,fit,decorations.pathreplacing}
\usepackage{pgfplots}
\usepgfplotslibrary{polar}
\pgfplotsset{compat=1.18}
\usepackage[edges]{forest}

% ============ 页眉页脚（报告标题和日期需在main.tex中覆盖） ============
\usepackage{fancyhdr}
\pagestyle{fancy}
\fancyhf{}
\fancyhead[L]{\scriptsize\color{steelgrey} \reporttitle}
\fancyhead[R]{\scriptsize\color{steelgrey} \reportdate}
\fancyfoot[C]{\small\color{steelgrey} \thepage}
\fancyfoot[R]{\tiny\color{steelgrey} \reportdisclaimer}
\renewcommand{\headrulewidth}{0.4pt}
\renewcommand{\footrulewidth}{0.2pt}

% ============ 章节格式 ============
\usepackage{titlesec}
\titleformat{\chapter}[block]
  {\normalfont\LARGE\bfseries\color{navy}}
  {\thechapter}{0.9em}{}
\titlespacing*{\chapter}{0pt}{0.4em}{1.0em}
\titleformat{\section}
  {\normalfont\large\bfseries\color{navy}}
  {\thesection}{1em}{}
\titleformat{\subsection}
  {\normalfont\normalsize\bfseries\color{steelgrey}}
  {\thesubsection}{1em}{}

% ============ 超链接 ============
\usepackage{hyperref}
\hypersetup{
  colorlinks=true,
  linkcolor=navy,
  urlcolor=navy,
  citecolor=navy,
}

% ============ 其他工具 ============
\usepackage{enumitem}
\setlist{nosep, leftmargin=2em}
\IfFileExists{fontawesome5.sty}{\usepackage{fontawesome5}}{}
% ============ Class B 全局兜底（xeCJK + 缺字 fallback 阻断）：
% 当 fontawesome5 因 Times New Roman 缺字 / \FA 命名空间失败时，
% 禁止输出 aExclamationTriangle / aCloud 裸文本，一律降级为 amssymb 原生符号。
% 注：amssymb 在 L112 已加载，此处不重复。所有 \def 在本段落。
\def\faExclamationTriangle{\ensuremath{\textcolor{riskred}{\blacktriangle}}\kern-0.25em}
\def\faCloud{\ensuremath{\clubsuit}\kern-0.1em}
\def\faYenSign{\ensuremath{\yen}}
\def\faIndustry{\ensuremath{\S}}
\def\faCogs{\ensuremath{\ast}}
\def\faFlagCheckered{\ensuremath{\blacktriangleright}}
\def\faMemory{\textbf{[RAM]}}
\def\faServer{\textbf{[DC]}}
\def\faGlobe{\ensuremath{\diamond}}
\def\faStar{\ensuremath{\bigstar}}
\def\faMicrochip{\textbf{[$\mu$P]}}
\def\faLayerGroup{\textbf{[LAY]}}
\def\faCubes{\ensuremath{\boxtimes}}
\usepackage{amssymb}
\usepackage{pifont}
\usepackage{tcolorbox}
\tcbuselibrary{skins,breakable}
\usepackage{caption}
\captionsetup{font=small, skip=4pt}
\usepackage{float}
\usepackage{needspace}
\setlength{\textfloatsep}{8pt plus 2pt minus 2pt}
\setlength{\floatsep}{6pt plus 2pt minus 2pt}
\setlength{\intextsep}{6pt plus 2pt minus 2pt}
\setlength{\abovecaptionskip}{4pt}
\setlength{\belowcaptionskip}{2pt}

% ============ 自定义命令 ============
\newcommand{\rmark}{\textcolor{riskred}{\ding{108}}}
\newcommand{\amark}{\textcolor{riskamber}{\ding{108}}}
\newcommand{\gmark}{\textcolor{riskgreen}{\ding{108}}}
\newcommand{\starfive}{$\bigstar\bigstar\bigstar\bigstar\bigstar$}
\newcommand{\starfour}{$\bigstar\bigstar\bigstar\bigstar$}
\newcommand{\starthree}{$\bigstar\bigstar\bigstar$}
% 资料来源注：使用楷体（STKaiti）而非斜体，避免中文无斜体字形导致的字体回退警告
% \mdseries 避免继承周围 \bfseries 导致 STKaiti 粗体缺失警告
% \sloppy + \hbadness 允许长文件名/URL 在必要时断行
% 局部重定义 \textbf → \heiti：STKaiti 无粗体，用黑体替代强调
\newcommand{\sourcenote}[1]{\par{\scriptsize\mdseries\color{mutedgrey}\kaishu\sloppy\hbadness=9999\def\textbf##1{{\heiti ##1}} 资料来源：#1}\vspace{0.2em}}
\newcommand{\qualitynote}[1]{\par{\scriptsize\mdseries\color{mutedgrey}\sloppy\def\textbf##1{{\heiti ##1}} 数据质量：#1}\vspace{0.2em}}
\newcommand{\evidencetag}[1]{%
  \tikz[baseline=(x.base)]\node[
    xshift=0pt,
    inner xsep=3.5pt,
    inner ysep=1.2pt,
    rounded corners=1.5pt,
    fill=panelgrey,
    draw=rulegrey,
    text=steelgrey,
    font=\scriptsize
  ] (x) {#1};%
}
\newcommand{\stance}[2]{\textbf{\textcolor{#1}{#2}}}

% 信息框
\newtcolorbox{keyinsight}[1][]{
  colback=lightgrey, colframe=navy,
  fonttitle=\bfseries, title={#1}, breakable,
  arc=2pt, boxrule=0.5pt, left=8pt, right=8pt, top=6pt, bottom=6pt
}
\newtcolorbox{riskbox}[1][]{
  colback=red!5, colframe=riskred,
  fonttitle=\bfseries, title={#1}, breakable,
  arc=2pt, boxrule=0.5pt, left=8pt, right=8pt, top=6pt, bottom=6pt
}
\newtcolorbox{dashboardbox}[1][]{
  enhanced, breakable,
  colback=white, colframe=navy,
  boxrule=0.7pt, arc=2pt,
  fonttitle=\bfseries\color{white},
  coltitle=white, colbacktitle=navy,
  title={#1},
  left=9pt, right=9pt, top=7pt, bottom=7pt
}
\newtcolorbox{houseviewbox}[1][]{
  enhanced, breakable,
  colback=lightgrey, colframe=deepnavy,
  borderline west={3pt}{0pt}{accentblue},
  boxrule=0.4pt, arc=2pt,
  fonttitle=\bfseries, title={#1},
  left=9pt, right=9pt, top=7pt, bottom=7pt
}
\newtcolorbox{sourcequalitybox}[1][]{
  enhanced, breakable,
  colback=yellow!6, colframe=riskamber,
  boxrule=0.45pt, arc=2pt,
  fonttitle=\bfseries, title={#1},
  left=9pt, right=9pt, top=7pt, bottom=7pt
}
% exhibitbox：表/图外框。left=7pt+right=7pt=14pt 内边距，before upper 内的 \linewidth 已扣除该内边距。
% 约定：\exhibitboxwidth = 真实内边距宽（tcolorbox 已扣好 padding 的 \linewidth）。
% 由于 tabularx/longtable 宽度参数的语义是「外部总宽（含两侧 \tabcolsep）」，内部列实际可用 = W - 2\tabcolsep。
% 所以 tabularx 直接传 \exhibitboxwidth（总宽 == 内边距宽），而其他用途（resizebox、tikz、p{...}）用 \exhibitwidth = \exhibitboxwidth - 2\tabcolsep。
\newlength{\exhibitboxwidth}
\newlength{\exhibitwidth}
\newtcolorbox{exhibitbox}[1][]{
  enhanced, breakable,
  colback=white, colframe=rulegrey,
  borderline north={1.2pt}{0pt}{navy},
  boxrule=0.35pt, arc=1pt,
  fonttitle=\bfseries\color{navy}, title={#1},
  left=7pt, right=7pt, top=6pt, bottom=6pt,
  before upper={\global\exhibitboxwidth=\linewidth\relax
    \global\exhibitwidth=\dimexpr\linewidth-2\tabcolsep\relax},
}
\newtcolorbox{disclosurebox}[1][]{
  enhanced, breakable,
  colback=lightgrey, colframe=steelgrey,
  boxrule=0.35pt, arc=1pt,
  fonttitle=\bfseries, title={#1},
  left=8pt, right=8pt, top=6pt, bottom=6pt
}

% ============ First-page dashboard cover ============
\newcommand{\astockcover}{%
\begin{titlepage}
\thispagestyle{empty}
\begin{tikzpicture}[remember picture, overlay]
  \fill[deepnavy] (current page.north west) rectangle ([yshift=-2.55cm]current page.north east);
  \fill[accentblue] ([yshift=-2.55cm]current page.north west) rectangle ([yshift=-2.66cm]current page.north east);
  \fill[deepnavy] (current page.south west) rectangle ([yshift=1.2cm]current page.south east);
\end{tikzpicture}
\vspace*{0.62cm}
{\color{white}\sffamily\bfseries\small \reportkicker}\\[0.18cm]
{\color{steelgrey}\sffamily\scriptsize \reportscope \quad|\quad \reporttype \quad|\quad \reportdate}\\[1.0cm]
{\fontsize{25}{30}\selectfont\bfseries\color{navy}\reporttitle}\\[0.25cm]
{\fontsize{13}{17}\selectfont\color{steelgrey}\reportsubtitle}\\[0.75cm]
\begin{dashboardbox}[决策摘要]
\begin{tabularx}{\textwidth}{L{3.0cm}X L{3.0cm}X}
\textbf{报告类型} & \reporttype & \textbf{数据截至} & \reportdatacutoff \\
\textbf{覆盖范围} & \reportscope & \textbf{作者} & \reportauthor \\
\end{tabularx}
\vspace{0.45em}
\begin{houseviewbox}[本机构观点]
\reporthouseview
\end{houseviewbox}
\begin{sourcequalitybox}[证据质量]
\reportquality
\end{sourcequalitybox}
\end{dashboardbox}
\vfill
{\scriptsize\color{steelgrey}\reportdisclaimer}
\end{titlepage}
}

\usepackage{amsmath}
\newcolumntype{P}{>{\raggedright\arraybackslash}X}
\hypersetup{pdfauthor={\reportauthor},pdftitle={\reporttitle}}
\fancyfoot[R]{}

\begin{document}
\astockcover

\begin{abstract}
本报告回答一个具体问题：固定20万元全额投资本金，如何在不把历史分红或某日价格误写为收益承诺的前提下，建立可复核的长期多资产风险预算？

核心结论：本模型明确指定CNY70,000的\texttt{511010}、CNY30,000的\texttt{510880}、条件性CNY30,000的\texttt{512890}和CNY20,000的\texttt{518880}；另将CNY40,000放入合格受保障流动性并保留CNY10,000交易现金。风险资产在资格包通过后按六个月计划分批建立；任一ETF的产品、价格/IOPV、折溢价、价差、流动性或组合穿透未通过时，该笔资金不交易并唯一留在受保障流动性。公募REITs和直接高股息股票初始均为0\%。长期年化3\%--6\%只作为五年以上的公式敏感性区间；实际现金仅以税后、扣费后到账额$D_0$记录。

该框架不提供订单、目标价、短线择时或保证收益。当前执行日资格包尚未运行，报告不表示任何ETF已具备实施资格。基准稳健核心在12个月共同承压规划情景下的算术损失为16.30\%，即固定本金约CNY32,600，目的在于检验风险承受与流动性；它不会自行触发账户实际回撤控制。
\end{abstract}

{\small
\setcounter{tocdepth}{0}
\tableofcontents
\par}
\clearpage

\chapter{投资结论}
\begin{houseviewbox}[核心结论]
固定CNY200,000采用ETF/资产层优先的基准稳健核心模型：20\%组合内受保障流动性、5\%交易现金、35\%高质量人民币利率债、30\%红利权益ETF和10\%非杠杆黄金ETF。直接股票和公募REITs初始配置均为0\%。截至本报告，ETF执行日资格包尚未运行；任何未通过或缺少数据的ETF额度均留在受保障流动性。
\end{houseviewbox}

\section{模型标的与条件性动作}

对本案固定CNY200,000本金，以下是明确的条件性模型目标：先建立防守层；仅在执行日资格包通过后建立利率债，并在六个月内分批建立红利权益和黄金。它们是风险预算中的指定标的、金额、条件和后备位置，不是个性化适当性结论、交易委托或以历史价格推导的数量。

\begin{exhibitbox}[固定CNY200,000的指定模型标的]
\small
\begin{tabularx}{\exhibitboxwidth}{L{2.4cm}L{3.0cm}L{1.7cm}P}
\toprule
\textbf{资产层} & \textbf{指定模型标的} & \textbf{目标金额} & \textbf{条件性模型动作} \\
\midrule
受保障流动性 & 合格受保障存款或期限匹配的直接政府债 & CNY40,000 & T0建立；不以基金或理财替代受保障层。 \\
交易现金 & 账户结算现金或银行活期；\texttt{511880}仅为可选补充 & CNY10,000 & T0保留，用于分批实施和再平衡；硬上限为CNY10,000。 \\
利率债 & \texttt{511010} 国泰上证5年期国债ETF & CNY70,000 & 仅在完整执行日资格包使$q=1$后于T0建立。 \\
红利权益A & \texttt{510880} 华泰柏瑞上证红利ETF & CNY30,000 & 仅在$q=1$后按六个月计划分批建立。 \\
红利权益B & \texttt{512890} 华泰柏瑞中证红利低波动ETF & CNY30,000 & 仅在$q=1$且与A的合并持仓穿透通过后分批建立。 \\
黄金 & \texttt{518880} 华安易富黄金ETF & CNY20,000 & 仅在$q=1$后按六个月计划分批建立；不使用期货或杠杆。 \\
REITs与直接股票 & 无初始标的 & CNY0 & 不买入；分别等待项目级或发行人级独立研究。 \\
\bottomrule
\end{tabularx}
\end{exhibitbox}

执行日资格包必须保存同一中国市场时段的来源、时间戳、产品文件、价格/IOPV、折溢价、买卖价差、流动性、账户和最小交易单位资料；对\texttt{512890}还须保存与\texttt{510880}合并后的持仓穿透。缺任一字段或未通过受控阈值时，相关$q=0$，不发生该笔交易。对应金额唯一留在合格受保障流动性；不会增加交易现金至CNY10,000以上，也不会自动转入\texttt{511010}、其他ETF、REITs或直接股票。当前资格包状态为未运行，因此本报告不表示任何ETF已通过。

\section{策略定位}

本案是一个长期、条件性的风险预算框架，服务于固定本金的流动性、分散和再平衡安排，不承诺某个分红率、票息、净值回报或现金提取额。

现金分配与总回报必须分开看待：ETF分配来自基金资产，除权日会影响净值；利率债ETF、货币工具和黄金ETF也不应被写成固定生活费来源。报告只记录实际税后、扣费后到账现金$D_0$，并按权重、上限和资格开关再投资。

\section{基准稳健核心版}

\begin{exhibitbox}[固定CNY200,000的风险预算]
\small
\begin{tabularx}{\exhibitboxwidth}{L{3.2cm}P P P}
\toprule
\textbf{资产层} & \textbf{权重} & \textbf{金额} & \textbf{角色与约束} \\
\midrule
组合内受保障流动性 & 20\% & CNY40,000 & 防守与用途匹配；仅合格存款或到期前直接政府债可归入该层，也承接未通过ETF核验的余额。 \\
交易现金 & 5\% & CNY10,000 & 分批实施与再平衡缓冲；不等同于存款保险，且不得超过CNY10,000。 \\
高质量人民币利率债 & 35\% & CNY70,000 & 利率与防守资产层；仍须接受久期、净值和流动性风险及执行日资格核验。 \\
红利权益ETF & 30\% & CNY60,000 & 条件性权益与红利因子暴露；\texttt{512890}未通过穿透时保持0。 \\
非杠杆黄金ETF & 10\% & CNY20,000 & 宏观风险分散；不计入收入预算。 \\
公募REITs & 0\% & CNY0 & 只有三只独立项目通过项目、估值和流动性审查时，才可从利率债桶转入最多5\%。 \\
直接高股息股票 & 0\% & CNY0 & 观察池，不是初始配置；未来仅可替换红利权益ETF暴露，合计不超过15\%。 \\
\midrule
\textbf{合计} & \textbf{100\%} & \textbf{CNY200,000} & \\
\bottomrule
\end{tabularx}
\end{exhibitbox}

\section{回报和损失的正确表述}

长期年化3\%--6\%是五年以上规划时的公式敏感性区间，不是“最可能收益”、年度收益率、产品分配率或收益保证。参数$r$仅表示税费、交易成本和基金费用之前的名义总回报算术输入；它不生成或估计实际现金$D_0$。基准稳健核心在12个月共同承压情景下的算术结果为-16.30\%，固定本金对应约CNY32,600；这用于检验风险承受与流动性，不是止损线、最大损失承诺或实际回撤触发。

\section{适用条件与直接股票边界}

{\small
\begin{itemize}
  \item 资金期限至少五年，且三年内没有确定的购房、教育、医疗、债务偿还或其他刚性用途；若用途改变，应另建资金计划。
  \item 模型标的须在执行日复核法律文件、费率、价格/IOPV、折溢价、价差、流动性、持仓重叠和账户适当性；本报告当前不含通过记录。
  \item 直接股票只有在独立完成价格、估值、派息覆盖、现金流、税务、行业集中度和流动性核验后，才能替换红利权益ETF暴露；合计上限15\%、单一发行人上限5\%。不使用杠杆、黄金期货或借款，分配下降也不自动触发加仓。
\end{itemize}
\par}


\chapter{红利权益观察池与边界}
\section{直接股票排除与启用边界}

本案当前策略不持有直接股票，也不持有期货。直接股票初始权重为0\%，不设发行人名单、目标价、历史收益率或现金分配假设；期货初始权重同为0\%，不作为替代配置工具。

\begin{exhibitbox}[当前边界：直接股票与期货均不进入初始策略]
\small
\begin{tabularx}{\exhibitboxwidth}{L{3.0cm}P P}
\toprule
\textbf{风险资产类别} & \textbf{初始权重} & \textbf{当前处理} \\
\midrule
直接股票 & 0\% & 不建立持仓、不形成标的排序，也不纳入现金、估值或执行结论。 \\
期货 & 0\% & 不使用杠杆或衍生品替代现货基金配置。 \\
\bottomrule
\end{tabularx}
\end{exhibitbox}

\section{排除的原因}

现行材料没有形成可独立复核的发行人级估值、现金流、税务、集中度和流动性证据包。因此，任何直接股票敞口都会把未完成的发行人判断带入组合；在这些判断完成前，零权重是策略约束，而不是等待选择的观察池。

\section{未来版本的启用门槛}

以下为未来经独立复核的策略修订所必须满足的边界，不构成当前启用授权：

\begin{enumerate}
  \item 对每一发行人完成独立、时点明确的估值审查，并说明估值结论的证据边界；
  \item 核验现金流、资本开支、分配覆盖和资产负债表约束，不以历史现金分配代替前瞻判断；
  \item 按适用账户和持有期完成税务审查，不将基金规则直接套用于直接股票；
  \item 穿透检验与现有权益基金的持仓重叠、行业／集团集中度及相关性；
  \item 完成可交易性、容量、停复牌、成本和退出路径审查，并在实施日重新核验；
  \item 取得独立复核和正式策略修订记录后，方可讨论是否启用。
\end{enumerate}

即使未来所有门槛通过，直接股票总额仍不得超过组合的15\%，单一发行人不得超过5\%。任何拟议敞口须受同等或更严格的集中度、流动性和治理约束；未通过、过期或缺失的任一门槛均意味着直接股票权重保持0\%。

\sourcenote{本章仅陈述当前策略边界。已退役的历史行情、分配和期货资料均不是本章输入，且不得支持目标、配置、估值、现金或执行结论。}


\chapter{组合构建与权重}
\section{组合结构}

基准稳健核心版按资产层而不是按单只股票构建。组合内受保障流动性与交易现金合计25\%，用来降低被动卖出风险；利率债为35\%，红利权益ETF为30\%，非杠杆黄金ETF为10\%。直接股票和公募REITs在初始向量中均为0\%。目标权重不等于当前已通过实施状态：截至本报告，ETF资格包未运行，未通过额度按规则留在受保障流动性。

\begin{exhibitbox}[基准稳健核心版]
\small
\begin{tabularx}{\exhibitboxwidth}{L{3.1cm}P P P}
\toprule
\textbf{资产层} & \textbf{目标权重} & \textbf{固定金额} & \textbf{资格与上限} \\
\midrule
组合内受保障流动性 & 20\% & CNY40,000 & 合格存款或到期前直接政府债；唯一承接未通过ETF资格检查的余额，安全上限100\%。 \\
交易现金 & 5\% & CNY10,000 & 硬上限CNY10,000；货币工具不自动等同于受保障存款，也不承接延后ETF资金。 \\
高质量人民币利率债 & 35\% & CNY70,000 & 指定\texttt{511010}；产品、久期和交易质量$q=1$后才可实施。 \\
红利权益ETF A & 15\% & CNY30,000 & 指定\texttt{510880}；产品与交易质量$q=1$后才可分批。 \\
红利权益ETF B & 15\% & CNY30,000 & 指定\texttt{512890}；另须通过合并行业和集团穿透，否则保持0。 \\
非杠杆黄金ETF & 10\% & CNY20,000 & 指定\texttt{518880}；非杠杆、价格/IOPV和交易质量$q=1$后才可分批，硬上限12\%。 \\
公募REITs & 0\% & CNY0 & 仅在三只独立项目同时通过时，从利率债转入最多5\%。 \\
直接高股息股票 & 0\% & CNY0 & 观察池；未来只能替换红利权益ETF，合计上限15\%。 \\
\midrule
\textbf{合计} & \textbf{100\%} & \textbf{CNY200,000} & \\
\bottomrule
\end{tabularx}
\end{exhibitbox}

\section{指定模型标的与资格路径}

执行日资格包是一个可复核的字段清单，而非概括性的“交易质量”标签。每只ETF需记录同一中国市场时段的产品文件、来源、时间戳、价格和IOPV/参考值、折溢价、最佳买卖价、价差、成交额或可执行深度、最小交易单位和账户可行性。模型控制要求绝对折溢价不高于1.00\%、买卖价差不高于中间价的0.50\%、单笔不高于记录的同日成交额或可执行深度的10\%。这些是风险控制规则，不是当前市场观察。

\begin{exhibitbox}[指定模型标的]
\small
\begin{tabularx}{\exhibitboxwidth}{L{2.4cm}L{3.1cm}L{1.7cm}P}
\toprule
\textbf{资产层} & \textbf{指定标的} & \textbf{目标金额} & \textbf{资格条件与失败时动作} \\
\midrule
高质量人民币利率债 & \texttt{511010} 国泰上证5年期国债ETF & CNY70,000 & 当前产品文件、久期、价格/IOPV、折溢价、价差和流动性均通过使$q=1$；否则不交易，CNY70,000留在受保障流动性。 \\
红利权益ETF A & \texttt{510880} 华泰柏瑞上证红利ETF & CNY30,000 & 当前产品文件、价格/IOPV、折溢价、价差和流动性均通过使$q=1$；失败批次留在受保障流动性。 \\
红利权益ETF B & \texttt{512890} 华泰柏瑞中证红利低波动ETF & CNY30,000 & 除交易资格外，合并行业不得超过总组合25\%，前十大集团不得超过权益桶35\%；不通过则保持0，CNY30,000留在受保障流动性。 \\
非杠杆黄金ETF & \texttt{518880} 华安易富黄金ETF & CNY20,000 & 核验非杠杆现货暴露、文件、IOPV、折溢价、价差和流动性；失败批次留在受保障流动性。 \\
交易现金工具 & 账户结算现金或银行活期；\texttt{511880}只作可选补充 & CNY10,000 & 交易现金上限为CNY10,000；若使用\texttt{511880}，另核验申赎/交易机制、价格偏离和非存款风险。 \\
\bottomrule
\end{tabularx}
\end{exhibitbox}

\section{六个月建仓动作}

每个金额是对CNY200,000的目标名义金额；场内最小交易单位和当日价格导致的零头向下取整，并留在受保障流动性。任何延后金额不为凑足目标而增加交易现金或换入相关资产。当前资格包未运行，故下表为通过资格后才会进入的计划顺序，而不是现在的交易指令。

\begin{exhibitbox}[六个月模型建仓动作]
\small
\begin{tabularx}{\exhibitboxwidth}{L{2.5cm}P P}
\toprule
\textbf{阶段} & \textbf{金额与目标} & \textbf{前提与失败处理} \\
\midrule
T0防守层 & CNY40,000受保障流动性、CNY10,000交易现金、CNY70,000的\texttt{511010} & 前两项按账户和存款边界核验；\texttt{511010}仅在$q=1$时建立，未通过的CNY70,000留在受保障流动性。 \\
第1--5个月 & 每月CNY5,000的\texttt{510880}、CNY5,000的合格\texttt{512890}、CNY3,333的\texttt{518880} & 各ETF分别$q=1$；\texttt{512890}还须穿透合格。任一失败时仅延后该笔，并将该笔留在受保障流动性。 \\
第6个月 & CNY5,000的\texttt{510880}、CNY5,000的合格\texttt{512890}、CNY3,335的\texttt{518880} & 以可交易最小单位向下取整；所有余量与未通过金额均留在受保障流动性，不产生补单。 \\
公募REITs或直接股票 & CNY0 & 不因等待分配、市场下跌或风险资产延后自动启用；须独立通过预先设定门槛。 \\
\bottomrule
\end{tabularx}
\end{exhibitbox}

\section{追加与再平衡}

每季度记录实际到账现金$D_0$、资产权重、产品事件、集中度、外部现金流和实际回撤$DD_t$；每半年才进行常规再平衡，优先使用新增资金和实际分配，避免把价格下跌自动解释为加仓信号。以下情形启动人工复核：受保障流动性、利率债或总权益相对目标偏离超过5个百分点；交易现金超过CNY10,000或在分批期内低于目标；黄金偏离超过2个百分点；任何资产超过硬上限或失去资格；或有完整账户记录的$DD_t\leq-12\%$。压力情景数值不代替实际$DD_t$。


\chapter{规划区间与现金分配口径}
\section{收益与现金分配的正确口径}

本报告不为个股、ETF、利率债、黄金或组合给出点状年化收益、固定票息或固定现金分配预测。资产管理产品不能承诺本金和回报；ETF分配也来自基金资产而非额外财富。以某一期历史分配率、某日价格或“股息率加增长率”直接合成组合收益，无法反映价格、估值、税费、费用、再投资时点和分配变化。

\section{长期规划敏感性，不是预测}

下表只回答“若长期名义总回报参数$r$落在不同规划输入，十年公式结果会怎样”，而不判断哪个参数最可能发生。$r$是扣除税费、交易成本和基金费用之前的假设总回报输入；不含年度追加、取现、再投资时点或路径，也不生成实际现金$D_0$。

\begin{exhibitbox}[五年以上规划参数]
\small
\begin{tabularx}{\exhibitboxwidth}{L{3cm}P P}
\toprule
\textbf{年化总回报参数} & \textbf{在本报告中的用途} & \textbf{明确不代表} \\
\midrule
3.0\% & 保守端的长期复利公式敏感性输入 & 下一年度收益、利息或股息承诺。 \\
4.5\% & 中间的长期复利公式敏感性输入 & “基准预测”或可提取现金。 \\
6.0\% & 较高端的长期复利公式敏感性输入 & 目标收益、估值上行或产品分配保证。 \\
\bottomrule
\end{tabularx}
\end{exhibitbox}

规划区间为年化3\%--6\%，适用期限为五年以上。它只用于比较期限、损失预算和复利算术，不替代执行日前的产品和市场复核。

\section{实际现金分配的记录方法}

季度复核时，定义

\[
D_0=\sum \text{过去12个月实际到账的利息、分配和股息（税后、扣费后）}.
\]

$D_0$ 是已经发生的现金记录，不外推为未来分配率。每一笔记录须包含到账日、来源资产层、税前金额、税额、费用、净额和结算或证据凭证。若实际分配被削减，使用 $D_h=D_0(1-h)$，其中 $h\in\{0,25\%,50\%,100\%\}$，仅评估可再投资现金减少对复利速度的影响；只有已有实际$D_0$记录时才计算该敏感性。任何分配为零的年份也不改变组合风险上限；资金可以留在组合内受保障流动性，不能以杠杆或追高补回。

\section{不纳入本模型的数字}

\begin{itemize}
  \item 未来一年的组合回报、股息率、债券ETF票息、货币工具收益或黄金收益；
  \item 每年可提取的固定现金、生活费覆盖额、税后收入或“收益来源占比”；
  \item 基于观察池的个股等权收益、资本利得假设或目标价格；
  \item 未经执行日核验的ETF价格、净值、折溢价、成交量、价差和交易成本。
\end{itemize}

\qualitynote{长期公式敏感性与实际现金台账彼此独立：$r$不预测$D_0$，$D_0$也不改变已披露的公式案例。}


\chapter{再投资与复利敏感性}
\section{复利敏感性}

复利可来自资产净值变化和实际收到后再投入的现金，但本章不把两者混成固定分红假设。下式仅使用名义总回报参数$r$作税费、交易成本和基金费用之前的公式演示；没有年度追加、取现、再投资时点、通胀或市场路径。它不记录、估计或承诺$D_0$。

\[
V_{10}=\text{CNY200,000}\times(1+r)^{10}.
\]

\begin{exhibitbox}[十年总回报公式敏感性]
\small
\begin{tabularx}{\exhibitboxwidth}{P P P P}
\toprule
\textbf{年化参数} & \textbf{初始本金} & \textbf{第十年公式终值} & \textbf{公式累计增值} \\
\midrule
3.0\% & CNY200,000 & CNY268,783 & CNY68,783 \\
4.5\% & CNY200,000 & CNY310,594 & CNY110,594 \\
6.0\% & CNY200,000 & CNY358,170 & CNY158,170 \\
\bottomrule
\end{tabularx}
\end{exhibitbox}

表中数值只由公式计算，不是对十年后资产规模的承诺，也不含交易成本、税费、基金费用、年度追加、实际分配、再投资时点、通胀或市场路径。

\section{实际分配的再投资规则}

\begin{enumerate}
  \item 每季记录$D_0$的到账日、税费、费用、资产来源、净额和结算凭证；不以公告金额、历史收益率或公式参数$r$代替到账记录。
  \item 只有当资产层的资格开关为1且未超过硬上限时，实际现金才可补向该资产层。
  \item 优先用实际现金和新增资金恢复目标权重；若无合格缺口，资金留在组合内受保障流动性。
  \item 若用户将现金或本金改作消费、教育、医疗、购房或债务偿还，本报告的“全部再投资”假设即失效，必须另建用途和税后现金流模型。
\end{enumerate}

\section{分配削减敏感性}

\begin{exhibitbox}[已实现现金的审慎重算]
\small
\begin{tabularx}{\exhibitboxwidth}{L{3cm}P P}
\toprule
\textbf{情景} & \textbf{计算} & \textbf{含义} \\
\midrule
不削减 & $D_0$ & 仅反映过去12个月实际净到账现金。 \\
削减25\% & $0.75D_0$ & 可再投资现金下降，非概率判断。 \\
削减50\% & $0.50D_0$ & 不以杠杆或卖出防守资产弥补。 \\
全部归零 & $0$ & 组合风险上限不变；只是没有可再投资现金。 \\
\bottomrule
\end{tabularx}
\end{exhibitbox}

\qualitynote{本案不设“第一年分红”“十年累计分红”或固定取现额；只有产生实际现金记录后，才可衡量$D_0$及其削减敏感性。}


\chapter{风险控制与压力测试}
\section{12个月共同承压规划测试}

以下冲击用于检验风险预算，不是历史回测、VaR、概率预测或最大损失承诺。每一资产层同时承受假设的不利变动；输入没有按历史频率或置信区间校准。交易成本、税费、跟踪误差、折溢价和再投资时点没有纳入。

\begin{exhibitbox}[基准稳健核心的共同承压算术]
\small
\begin{tabularx}{\exhibitboxwidth}{L{3cm}P P P}
\toprule
\textbf{资产层} & \textbf{权重} & \textbf{共同承压假设} & \textbf{加权贡献} \\
\midrule
组合内受保障流动性 & 20\% & 0\% & 0.00\% \\
交易现金 & 5\% & -1\% & -0.05\% \\
高质量人民币利率债 & 35\% & -5\% & -1.75\% \\
红利权益ETF & 30\% & -45\% & -13.50\% \\
非杠杆黄金ETF & 10\% & -10\% & -1.00\% \\
\midrule
\textbf{组合} & \textbf{100\%} &  & \textbf{-16.30\%} \\
\bottomrule
\end{tabularx}
\end{exhibitbox}

固定CNY200,000本金下，-16.30\%对应约CNY32,600。这一数字是共同承压情景的算术输出，用于讨论损失承受和流动性；它不会自行触发实际回撤复核，也不说明实际损失会停在该数值。

\section{实际回撤线与响应}

实际回撤必须以可追溯的账户价值记录测量：

\[
DD_t=\frac{\text{经外部现金流调整后的当期组合市值}}{\text{经外部现金流调整后的已记录高水位}}-1.
\]

每次记录至少应保留复核日期、账户或托管来源、当期市值、外部申购/取现、调整后高水位及数据质量。若这些字段缺失，不宣称已触发回撤线；情景表也不能替代该记录。

\begin{exhibitbox}[控制线]
\small
\begin{tabularx}{\exhibitboxwidth}{L{3cm}P P}
\toprule
\textbf{触发} & \textbf{阈值} & \textbf{研究响应} \\
\midrule
常规测量 & 每季度 & 记录权重、实际净分配、产品文件变化、折溢价、集中度和实际$DD_t$。 \\
再平衡复核 & 受保障流动性、利率债或总权益偏离目标5个百分点；交易现金超过CNY10,000或在分批期内低于目标；黄金偏离2个百分点 & 优先用新增资金和实际现金修复；不把下跌自动解释为买入信号。 \\
实际损失复核线 & 已记录$DD_t\leq-12\%$ & 暂停自动新增风险资产，重新核验资产层资格和用途匹配。 \\
实际风险处理线 & 已记录$DD_t\leq-18\%$ & 先检查流动性、杠杆、产品异常和权重，再由投资者作出非自动化决定。 \\
直接股票激活 & 合计不超过15\%，单一发行人不超过5\% & 只有发行人估值、现金流、税务、行业集中度和流动性审查全部通过。 \\
\bottomrule
\end{tabularx}
\end{exhibitbox}

\section{主要风险}

\subsection{利率、净值与流动性风险}

利率债ETF并非存款，利率上行时可发生净值下跌；货币ETF也不自动进入存款保险边界。任何产品在实施日前都需核验期限、费率、交易机制、折溢价和流动性。不能用“国债”或“货币”名称替代产品风险审查。

\subsection{权益因子和集中度风险}

红利权益ETF可以同时暴露于金融、周期和高分配行业。两只ETF并不天然分散，必须穿透看共同持仓、集团和行业权重。历史分配率下降、盈利变化或估值波动都可能使总回报和实际现金分配低于过去观察值。

\subsection{黄金与REITs风险}

黄金现货ETF不产生固定利息，且可能在共同承压中下跌。公募REITs有项目、估值和流动性风险；在三只独立项目未完成审查前，正确权重是0\%，而不是为了追求分配而提前买入。

\section{反方观点}

\textbf{质疑}：30\%红利权益ETF和35\%利率债仍可能同时下跌，3\%--6\%规划区间是否过于乐观？

\textbf{回应}：区间不是预测，也不用于证明策略“安全”。本案把共同承压、流动性缓冲、权重硬上限和实际回撤记录放在收益讨论之前；如果投资者无法接受约CNY32,600的规划压力损失，应选择更保守的向量或重新定义资金用途。


\chapter{实施研究计划与再平衡}
\section{条件性实施研究计划}

本章给出条件性模型仓位的建立顺序，不构成交易执行授权。未补齐原始证据、未取得实施日产品和市场数据并完成独立复核前，不得据此提交交易委托或确认价格、数量、通道或适当性。没有完整执行日资格包时，所有ETF资格开关均为$q=0$；各笔目标、条件及失败后的资金位置仍可复核。

\begin{exhibitbox}[六个月模型实施动作]
\footnotesize
\begin{tabularx}{\exhibitboxwidth}{L{2.6cm}P P}
\toprule
\textbf{阶段} & \textbf{模型金额} & \textbf{应完成的检查与失败处理} \\
\midrule
启动防守层 & CNY40,000受保障流动性、CNY10,000交易现金和CNY70,000的\texttt{511010} & 验证资金用途、存款保险边界、账户和产品资料；\texttt{511010}还须通过当前价格/IOPV、折溢价、价差和流动性检查。未通过的ETF金额留在受保障流动性。 \\
第1--5个月 & 每月CNY5,000的\texttt{510880}、CNY5,000的合格\texttt{512890}和CNY3,333的\texttt{518880} & 每笔记录同一市场时段的来源、时间戳、价格/IOPV、价差和流动性；\texttt{512890}另需合并持仓穿透。任一失败只延后该笔，并留在受保障流动性。 \\
第6个月 & CNY5,000的\texttt{510880}、CNY5,000的合格\texttt{512890}和CNY3,335的\texttt{518880} & 以最小交易单位向下取整；剩余与未通过金额留在受保障流动性，交易现金仍不得超过CNY10,000。 \\
REITs与直接股票 & CNY0 & 未通过独立激活门槛时保持0，不以市场下跌、等待分配或ETF延后绕过资格开关。 \\
\bottomrule
\end{tabularx}
\end{exhibitbox}

金额为风险预算的算术映射，唯一后备位置是受保障流动性：核验失败不会增加\texttt{511010}、替换ETF或令交易现金超过CNY10,000。延后额度只可在下一计划分批或再平衡日重取完整资格包后实施，不补单。

\section{实施日前的资格检查}

{\footnotesize
\begin{enumerate}
  \item 确认固定CNY200,000确属五年以上可承受波动的本金，且没有改变为刚性用途。
  \item 对每个候选ETF保存当前产品文件、费用、来源、时间戳、交易状态、价格/IOPV、折溢价、最佳买卖价、价差、同日成交额或可执行深度、账户资格、最小交易单位和现金可行性。
  \item ETF仅在绝对折溢价不高于1.00\%、买卖价差不高于中间价0.50\%、单笔不高于记录同日成交额或可执行深度10\%时$q=1$；缺字段或未通过即$q=0$。
  \item 对两只红利权益ETF做带日期的合并持仓穿透，确认行业不超过总组合25\%，前十大集团不超过权益桶35\%；\texttt{512890}未通过时保持0。
  \item 核验利率债久期与用途匹配且不把债券/货币ETF写成保本存款；黄金只可为非杠杆现货ETF；公募REITs须有至少三只独立项目通过现金流、债务、估值溢价、流动性和最小交易单位审查，才可从利率债转入最多5\%。
\end{enumerate}
\par}

\section{再平衡与记录}

{\footnotesize
\begin{itemize}
  \item \textbf{季度}：记录实际到账的税后、扣费后现金$D_0$、资产权重、产品事件、资格开关、外部现金流和实际$DD_t$；$D_0$须有到账日、来源、税费、净额和凭证。
  \item \textbf{半年}：优先用新增资金和实际现金分配恢复权重；受保障流动性、利率债或总权益偏离5个百分点、交易现金超过CNY10,000或分批期内低于目标、黄金偏离2个百分点时复核。
  \item \textbf{实际回撤与年度}：已记录$DD_t\leq-12\%$时停止自动新增风险暴露，$DD_t\leq-18\%$时先复核流动性、杠杆、产品异常和资金用途；情景损失不替代$DD_t$。年度重读产品文件和税务规则，更新现金、穿透和观察池；未独立研究时直接股票仍为0\%。
\end{itemize}
\par}

\qualitynote{本文件为条件性研究草案。实施日资格包、个人账户适当性、税务和任何交易决定须由投资者及其受托专业人士独立确认。}


\appendix
\begingroup
\small
\chapter{数据来源、证据边界与假设}
\section{来源登记与本地归档状态}

本案当前读者可用的来源池仅包含15项完整活跃材料：S01--S13及S15--S16。它们的原件已本地归档；来源权威性、资料时效和实施日资格仍须分别复核。

\begin{longtable}{L{2.0cm}L{4.1cm}L{9.1cm}}
\toprule
\textbf{活跃来源} & \textbf{来源类别} & \textbf{在本案中的有限用途} \\
\midrule
S01--S07 & 税务、监管及存款保险一级材料 & 支持税务、基金风险承担、存款保护、资管产品非保本及个人养老金等制度边界；不构成个体税务意见。 \\
S08 & 基金一级产品文件 & 支持利率债基金的产品身份及利率风险边界；不证明当前价格、流动性或实施资格。 \\
S09--S13 & 交易所及基金一级产品／投资者教育材料 & 支持权益、黄金及货币市场基金的产品身份、分配机制与风险边界；不证明当前收益、净值、折溢价或流动性。 \\
S15 & 交易所一级材料 & 支持公募REITs的结构与风险边界；不构成项目选择或估值结论。 \\
S16 & 指数公司一级材料 & 支持红利低波动暴露的集中度提醒；不替代实施日前的完整穿透审查。 \\
\bottomrule
\end{longtable}

\paragraph{退役的历史资料。}
S14、M01及D01--D08仅保留为审计历史，均标记为非活跃范围并已从当前范围退役，不是当前读者输入。它们不得支持任何目标、配置、估值、现金或执行结论；其未取得或不完整的历史原始字节也未被声称为已恢复。

\section{模型假设与边界}

\begin{itemize}
  \item 固定CNY200,000是全额可投资本金；组合内受保障流动性和交易现金属于该本金，不另从家庭资产负债表扣减。
  \item 基准稳健核心权重为20\%受保障流动性、5\%交易现金、35\%高质量人民币利率债、30\%红利权益ETF和10\%非杠杆黄金ETF；REITs、直接股票和期货初始均为0\%。
  \item 指定模型目标为CNY70,000的\texttt{511010}、CNY30,000的\texttt{510880}、条件性CNY30,000的\texttt{512890}和CNY20,000的\texttt{518880}。当前资格包未运行；任一ETF未通过时，相关金额留在受保障流动性，交易现金不超过CNY10,000。
  \item 年化3\%--6\%仅为五年以上、税费和费用之前的名义总回报公式敏感性，不是收益预测、固定现金分配、目标收益或个股／ETF估值结论；参数$r$不生成实际现金$D_0$。
  \item 所有现金分配以实际税后、扣费后到账额$D_0$记录并默认再投资；不预测第一年现金、十年累计分配或可消费金额。
  \item 压力测试为12个月的共同冲击算术，不是历史回测、VaR、概率估计或最大损失保证；只有账户记录计算的实际$DD_t$才可触发回撤复核。
\end{itemize}

\section{数据质量与草案状态}

活跃来源池只覆盖规则和产品边界，并不生成实施日资格。当前没有价格、IOPV、折溢价、价差、流动性、账户资格或合并持仓穿透记录；因此所有ETF资格开关仍为$q=0$，本文件仅为条件性研究框架。


\medskip
\begin{disclosurebox}[重要披露]
\footnotesize
\reportdisclaimer\par
直接股票初始为0，未有任何个股名称、价格或历史股息率进入配置、收益或现金流结论。实施日核验、账户适当性和独立复核完成前，本文件不得作为交易依据、适当性结论或收益／本金承诺材料。
\end{disclosurebox}
\endgroup

\end{document}
